\documentclass[11pt]{article}

\usepackage[margin=1in]{geometry}
\usepackage[utf8]{inputenc}
\usepackage[T1]{fontenc}
\usepackage{lmodern}
\usepackage{amsmath,amssymb,amsthm}
\usepackage{booktabs,array}
\usepackage{microtype}
\usepackage[numbers,sort&compress]{natbib}
\usepackage{xcolor}
\usepackage{etoolbox}
\usepackage{graphicx}
\usepackage[font=small,labelfont=bf]{caption}
\AtBeginEnvironment{table}{\small\setlength{\tabcolsep}{3pt}}
\usepackage[colorlinks=true,linkcolor=blue!50!black,citecolor=blue!50!black,urlcolor=blue!50!black]{hyperref}

% \pn, a key and its value: a number regenerated from the data (indrajala-studies'
% scripts/paper/certificate_note_tables.py check and fill)
\newcommand{\pn}[2]{#2}
\newcommand{\draftnote}[1]{\textcolor{red!70!black}{[\textbf{Draft:} #1]}}
% The preprint's DOI, reserved on its Zenodo draft before upload; empty, the title page shows a draft note
\newcommand{\notedoi}{}
% The addendum's DOI, reserved on its Zenodo draft; empty, the text shows a draft note
\newcommand{\addendumdoi}{}

\newtheorem{theorem}{Theorem}
\newtheorem{proposition}[theorem]{Proposition}
\theoremstyle{definition}
\newtheorem{definition}[theorem]{Definition}

\title{Proving a Trained Network Is a Finite-State Machine: an Exact Certificate, Its Witness and
Its Numerics\\[0.5em]
\large Research note, draft, November 2026\\[0.2em]
\ifx\notedoi\empty\draftnote{the DOI}\else\href{https://doi.org/\notedoi}{doi:\notedoi}\fi}

\author{
  David Barkhuizen\\
  Independent researcher
}

\date{}

\begin{document}

\maketitle

\begin{abstract}
A network trained on a task that a finite-state machine solves can pass every test and still be
wrong: tests sample, and a machine extracted from the network is not the network. This note decides
the question exactly, for networks built to allow it. When the only channel between a recurrent
network's steps is a finite code, each state replaced exactly by its nearest codebook row, the
network \emph{is} a Mealy machine, and its equivalence with a target machine is decided by enumerating
every reachable code under every input window. That core has a recent precedent. This note states it
as a method for steps that read the state, with its proof, near ties branched, a witness that an
independent checker of a few hundred lines verifies by re-running every entry through the forward pass,
and sound bounds on its floating-point decisions. We validate it on existing networks, training none.
All \pn{recert-certified}{1{,}077} certified tables of \pn{recert-certified-networks}{674} networks were
valid and all \pn{recert-uncertified}{11} uncertified ones invalid. A modelled sampled test of 100{,}000
words a length misses a planted fault at a rarely visited entry with chance up to
\pn{fault-division-worst}{0.9998} on modular division. Interval and affine bounds show that every entry's
decisions in the \pn{bound-tables}{613} tables they cover hold for the network's real-valued function, not
only for its floating-point evaluation. Without the target, the proof needs a bound on the target's size,
and we say where that holds.
\end{abstract}

\section{The problem, and why it matters}

\paragraph{The problem.} Train a recurrent network on a task that a finite-state machine solves:
adding numbers column by column, say, with the carry as the state. Test it on inputs far longer than
any it trained on. It passes. Is it right on every input, at every length?

Tests cannot say. In an earlier study, two addition networks passed every sampled test at every
length to 64 columns and are wrong at twelve columns. One of them adds $1{,}559 + 6{,}329 + 5{,}231$
to $19{,}680$. The other also passed all 328{,}000 of its long-test cases, to 1{,}024 columns
\citep[Appendix B]{barkhuizen-2026-learning-transducers-with-proof}. The problem this note solves is
this one: given a recurrent network built so that only a finite code crosses its steps, and a
finite-state machine, decide exactly whether they compute the same function on every input of every
length the network can read.

\paragraph{Why it matters.} Three kinds of work need that answer and mostly estimate it.
\begin{itemize}
\item \emph{Claims that a network learned a rule.} Studies of state tracking, length generalisation
and automaton extraction measure success by sampled agreement, which overstates it by an amount they
cannot see. An exact verdict gives them a ground truth.
\item \emph{A network standing in for a program.} A network is useful beyond its training lengths
only if it is right there. For a program one can ask for a proof; for a network one usually cannot.
\item \emph{Finding an unknown machine from its behaviour.} A learner needs to know when it has
finished. A proof gives it a stopping rule, and turns a success rate into a cost: a certified run has
an exact answer and a random running time.
\end{itemize}

\paragraph{The price.} The method does not apply to an existing network with a continuous state. It
asks for a design: make the recurrent interface finite. Without one the question can be out of reach
altogether: for general weighted ReLU recurrent networks, equivalence with an automaton is undecidable
\citep{marzouk-2020-distance-equivalence-fsm-rnn}. Given a finite interface, equivalence becomes decidable by
enumeration, and cheaply so for the networks here: at most \pn{recert-largest-entries}{5{,}314{,}410}
window evaluations, in minutes.

\paragraph{What is in the note.} Most existing evidence is of two kinds. Sampled tests can show that a
network is wrong, never that it is right. Extraction builds a surrogate machine from the network, by
quantising its trained state
\citep{giles-1992-learning-extracting-fsa-second-order-rnn,omlin-1996-extraction-rules-discrete-time-rnn},
by L$^\ast$ with the network as teacher \citep{weiss-2018-extracting-automata-queries-counterexamples}
or with sampled validation \citep{mayr-2018-regular-inference-artificial-neural-networks}; proving the
surrogate says nothing exact about the network, whose behaviour on unvisited states can differ. A third
kind is a decision. Its nearest precedent is \citet{brandoit-2026-length-independent-state-tracking}, who
read a trained network's transition tables and check them against the target by an exhaustive
exploration of the reachable codes, for steps whose tables do not depend on the state. This note states
the method for steps that read the state, under stated conditions (Section~\ref{sec:setting}), with its
proof (Section~\ref{sec:procedure}), what its numerics assume and how far sound bounds remove the
assumption (Section~\ref{sec:numerics}), the witness that makes it a certificate in the sense of
\citet{mcconnell-2011-certifying-algorithms} (Section~\ref{sec:witness}), what remains without the target
(Section~\ref{sec:target-free}), its uses as an instrument (Section~\ref{sec:instrument}), and its
validation (Section~\ref{sec:validation}). The method was the instrument behind
\citet{barkhuizen-2026-learning-transducers-with-proof}, where it has one paragraph of the method and
parts of the appendices; here it is the subject. The results cited validate it; they are not the point.

\section{The setting}
\label{sec:setting}

A \emph{Mealy machine} $T = (S, \Sigma, \Gamma, s_0, \tau, \lambda)$ has finite states $S$, input and
output alphabets $\Sigma$ and $\Gamma$, a start $s_0$, a transition $\tau : S \times \Sigma \to S$ and an
output $\lambda : S \times \Sigma \to \Gamma$. It maps every word to a word of the same length. We take
targets minimal (no two states equivalent) and every state reachable from $s_0$.

The \emph{window-recurrent network} reads a word $k$ tokens at a time. A window's step is a function
$F(c, x) = (y, h)$ of the state $c \in \mathbb{R}^d$ handed on by the previous window and the window's
$k$ tokens $x$: $y$ is the $k$ output distributions and $h$ the next state before it is snapped. The
state is \emph{factored}: $G$ factors, factor $g$ with a codebook $E_g$ of $m_g$ rows, and the snap
$\sigma$ replaces each factor's block of $h$ by its nearest codebook row. A \emph{joint code} is a tuple
$q = (q_1, \dots, q_G)$ with $q_g \in \{1, \dots, m_g\}$, and $E(q)$ its concatenated rows. In the
networks studied here the step is a small transformer over the window's tokens, with the state added
to its tokens, an attention pooling per factor and a GRU cell \citep{cho-2014-rnn-encoder-decoder-gru};
nothing below depends on that choice beyond the conditions that follow. Some networks read a window
through fixed query tokens, one input column followed by the queries, and answer in one output slot
(\pn{recert-adapter-tables}{452} of the tables here). The step is then the wrapped function
$F'(c, x)$ of that column, and the conditions apply to $F'$; so ``the machine is the network'' holds
under that stated reading, which the witness records.

\begin{definition}[the conditions]
\leavevmode
\begin{enumerate}
\item[C1] \emph{A finite cross-step interface:} the only value passed from one window to the next is
the snapped state, a joint code.
\item[C2] \emph{An exact snap:} the state handed on is exactly $E(q)$, bit for bit, so a window's
computation from code $q$ is the same function whatever the history that reached $q$.
\item[C3] \emph{No other channel:} a window's step reads only its tokens and the state; no position,
counter or cache crosses windows.
\item[C4] \emph{Deterministic evaluation:} $F$ is evaluated by a fixed program; the same inputs give
the same bits.
\end{enumerate}
\end{definition}

Under C1--C3 the network is the machine $N$ whose states are joint codes, with
\[
\delta_N(q, x) = \sigma\big(h(E(q), x)\big), \qquad \lambda_N(q, x) = \operatorname{argmax} y(E(q), x),
\]
over the alphabet $\Sigma^k$ of windows, started at the snapped initial state $q_0 = \sigma(h_0)$, a
constant, with $\sigma$ the Euclidean nearest code. C4 makes $\delta_N$ and $\lambda_N$ functions:
without it, the same entry could be read two ways. The implementation used here satisfies C2 by
construction for a loaded network, and a test checks it bitwise on small randomised networks: in
full-sequence inference the state handed on equals $E(q)$, and each window's outputs equal those computed
from $E(q)$ alone. C4 is checked on one platform. The earlier study retrained four networks on x86-64 and
aarch64 and found the step logs and the networks' SHA-256 identical
\citep{barkhuizen-2026-learning-transducers-with-proof}, which supports it across the two without
comparing certificate tables.

\paragraph{Words.} The machine reads whole windows. A word that does not fill its windows reaches the
network with leading zero columns, so a certificate speaks to it only if the zero token keeps the
target at its start with label zero. That holds for addition, the only family here at $k > 1$; at $k = 1$
every word is whole windows. We take $\Sigma^k$-reachability as what counts: at $k > 1$ a target's state
may be reachable by tokens and not by whole windows, and then no certificate needs it.

This is what separates the network from a surrogate. An extracted machine is a model of the network;
here the machine \emph{is} the network, read off without approximation, so a statement about it is a
statement about every input the network can be given.

\section{The procedure}
\label{sec:procedure}

\paragraph{Enumeration.} From $q_0$, breadth first, every window $x \in \Sigma^k$ is run from every
joint code reached, recording $D(q, x) = \lambda_N(q, x)$ and $N(q, x) = \delta_N(q, x)$, until no new
code appears. The reachable set $Q$ is finite (at most $\prod_g m_g$), and the search costs
$|Q| \cdot |\Sigma|^k$ window evaluations, batched.

\paragraph{The map to the target.} A map $\varphi : Q \to S$ is fixed as the search goes:
$\varphi(q_0) = s_0$, and a code first reached from $q$ by window $x$ gets
$\tau^\ast(\varphi(q), x)$, the target's state after $x$. An entry $(q, x)$ \emph{passes} when
\[
D(q, x) = \lambda^\ast(\varphi(q), x) \quad\text{and}\quad \varphi(N(q, x)) = \tau^\ast(\varphi(q), x),
\]
where $\lambda^\ast$ and $\tau^\ast$ extend the target's output and transition to windows. The network
is \emph{certified} when every entry passes and each label leads the next by a margin, with near ties
in the snap branched (Section~\ref{sec:numerics}).

\begin{theorem}[soundness and completeness]
\label{thm:exact}
For the network as evaluated (C4):
\begin{enumerate}
\item[(a)] \emph{Soundness.} If the network is certified, it computes the target's outputs on every
word of whole windows from the start, of every length.
\item[(b)] \emph{Completeness.} Read the decisions as evaluated, with no margin asked and no branch
taken. If some entry fails, the network and the target differ on some word. Let $(q, x)$ be the first
failure in breadth-first order, and $u$ the path by which $q$ was first reached. If a label is wrong,
$ux$ is such a word, and no word of fewer windows is one. If the next state is wrong, let $u'$ be
the path by which $N(q, x)$ was first reached, and $v$ a word telling the target's two states apart:
the network and the target differ on $uxv$ or on $u'v$.
\item[(c)] \emph{Refusals.} A failure of a margin or of a branch alone shows no word on which the
network is wrong under its own evaluation: the certificate has refused, not disproved. If no entry fails
with no branch taken, the network as evaluated is right, whatever the margins.
\end{enumerate}
\end{theorem}

\begin{proof}
(a) $\varphi$ is a homomorphism from $N$'s reachable part to $T$ over windows. By induction on the
number of windows, after a word $u$ the network is in a code $q$ with $\varphi(q)$ the target's state
after $u$, and the next window's labels are the target's from that state.

(b) Every entry before $(q, x)$ passes, so by the same induction, along any path through those
entries the network's labels are right and $\varphi$ gives the target's state. A wrong label therefore
makes $ux$ a word on which they differ. A word of fewer windows uses only entries of smaller depth,
all before $(q, x)$, and those pass, so it is right. Now let the next state be wrong: $q' = N(q, x)$
was first reached by $u'$, so the target is in $\varphi(q')$ after $u'$ and in a different state
$\tau^\ast(\varphi(q), x)$ after $ux$, while the network is in $q'$ after both. The target is minimal,
so a word $v$, padded to whole windows, makes the two target states' outputs differ. From $q'$ the
network gives one output on $v$, so it differs from the target after $u'$ or after $ux$.

(c) The margins and branches ask more than the evaluated decisions do, and (a)'s induction needs only
those decisions.
\end{proof}

With no branch taken, the search reports the first failing entry and, for a wrong state, the word $uxv$;
when that word is right, $u'v$ is the one that is wrong. Where the first failure is a wrong state, $uxv$
need not be a shortest counterexample, since a shorter wrong label may lie later at the same depth or the
next. With branching on, none of this is reliable: $\varphi$ can be fixed through a branch, so a reported
failure, its word and its counts may describe no real evaluation. A failure in a tie-branched run is
to be read through the strict run.

\paragraph{Minimisation.} The machine read off is minimised by Moore's refinement
\citep{moore-1956-gedanken-experiments-sequential-machines}; a certified network minimises to the size
of the target's part reachable by whole windows, and the codes that merge are the network's redundancy.
The tasks' alignment (leading zero columns) and final window are windows like any other, so under the
premise above a certificate covers every length and alignment the task allows.

\paragraph{Factored states.} With $G > 1$ factors the states are tuples and nothing above changes;
the search, $\varphi$ and the minimisation are over tuples. Each factor's \emph{readout} is the partition
of the target's states its codes induce over $Q$: whether the factor's code is a function of the
target's state, and which congruence of the target it holds. These readouts are exact, being read
off the certified machine, not estimated.

\section{Numerics}
\label{sec:numerics}

The proof is about the network as evaluated. Two decisions in each entry depend on floating-point
values: the label (an argmax) and the snap (an argmin of distances). A label is accepted only with a
lead of $10^{-6}$ in log-probability over the next; a near tie in a label fails the certificate. A
snap is a \emph{near tie} when another code's distance is within a ratio $\rho = 1.05$ of the nearest's.
Rounding elsewhere (another batch, another summation order) could snap such a state the other way,
so the search \emph{branches}: every code within the ratio is explored as a next state and checked
like the snapped one. A \emph{tie-branched certificate}\footnote{Called ``robustly certified'' in
\citet{barkhuizen-2026-learning-transducers-with-proof}; renamed here, since \emph{certified
robustness} means robustness to input perturbations in machine learning.} proves the network right
for every resolution of every near tie; a \emph{strict} one ($\rho = 1$) for its own snaps only. A
near tie that branches into a state the minimised machine tells apart from the snap's is a
\emph{split tie}: the machine's own nondeterminism, which can be counted without the target.

The margins are chosen constants. They are adequate if rounding moves no logit or distance by more
than they allow; that is an assumption, and \citet{jia-2021-exploiting-verified-nn-floating-point}
show why such assumptions should be stated and, where possible, checked. We check it with two sound
bounds, asking of each recorded decision: does the network's exact real-valued function make the same
one?

\paragraph{The interval bound.} Each entry's window is evaluated again in midpoint-radius interval
arithmetic \citep{moore-1966-interval-analysis,rump-1999-fast-parallel-interval-arithmetic} from its exact inputs (the window's tokens and the
code's rows). A sum of $n$ floats carries the order-free bound $\gamma_n = nu/(1-nu)$ times the sum of
its terms' magnitudes, $u = 2^{-53}$ \citep[Chapters~3--4]{higham-2002-accuracy-stability-numerical-algorithms};
every other elementary result is inflated by $u$ of its magnitude, with an absolute floor for underflow;
and $\exp$, $\tanh$ and the logistic function are assumed within a relative error of
\pn{bound-ulps}{4}$u$, an assumption on the maths library (square roots and reciprocals get the same
charge, though IEEE rounds them correctly). Softmax is enclosed as $p_i = 1/\sum_j \exp(z_j - z_i)$, so that no term appears
twice. The intervals then enclose the exact real value of every activation. A label is \emph{sound}
when its logit's interval lies above every other class's; a snap is sound when the nearest code's
squared-distance interval lies below every other's, and \emph{covered} when the codes whose intervals
overlap the nearest's are among the branches the certificate took. Intervals lose every correlation,
so across an attention block they widen by orders of magnitude, and on deep windows this bound decides
little.

\paragraph{The affine bound.} Affine arithmetic \citep{de-figueiredo-2004-affine-arithmetic-concepts-applications}
keeps the correlations. Each value is a form $c + \sum_s A_s e_s + [-r, r]$ with every $e_s \in [-1, 1]$
and the noise symbols $e_s$ shared across the entry's values. A linear map carries the coefficients
through exactly, so errors that cancel downstream cancel here. A product keeps both first-order terms
and bounds the second-order one by the radii's product. A nonlinear function is linearised at the
centre, with the remainder bounded through its second derivative over the input's range, and ReLU
across zero by the usual chord \citep{singh-2018-fast-effective-robustness-certification}. After each layer every remainder becomes a fresh symbol. Rounding is
charged as for intervals, along with the rounding of the bound's own arithmetic. Each entry keeps at
most \pn{affine-cap}{1{,}024} symbols, the heaviest; every dropped coefficient's magnitude joins its
value's remainder, so the form stays an enclosure and only loosens. Decisions are read on differences,
which keep the correlations: a label's lead over each other class, and for a snap to $c$ and another
code $e$, the difference $2\,h \cdot (E_c - E_e) + |E_e|^2 - |E_c|^2$ of their squared distances.

A decision that either bound finds sound is one the exact real-valued function makes too. Since the
recorded decisions came from the floating-point evaluation, the floating-point and exact machines then
agree on it. We rest no claim on other evaluations, though the error model is written to cover any
summation order.

Three things stand behind that sentence short of a proof. The bounds run on round-to-nearest with
inflated radii, not directed rounding, and neither program is formally verified: they extend the trusted
base (Section~\ref{sec:witness}) for this claim alone. The rounding of the radii themselves is not
charged exactly: the fixed inflation can fall short of a long sum's error, and a codebook row's computed
norm is taken as exact. For the exact function the charges for the network's own rounding are surplus and
exceed these, and in a spot check of the four affine tables with the tightest snaps (not in the tables
script), each difference's centre was more than $10^4$ times its radius. And when the cube of a reciprocal's
input overflows, the affine program drops that reciprocal's second-order charge; the value is then below
about $2 \times 10^{-103}$, far below the rounding charged downstream, and the logs record the event.

\section{The witness}
\label{sec:witness}

A certifying algorithm returns, with its output, a witness that a simple checker can verify
\citep{mcconnell-2011-certifying-algorithms}; translation validation checks each run of a compiler rather
than the compiler \citep{pnueli-1998-translation-validation}; and neural-network verifiers now produce
proofs for a separate checker
\citep{isac-2022-neural-network-verification-proof-production,desmartin-2024-certified-proof-checker-imandra}.
The certificate's witness is the table it read: every reachable code's labels, snapped next code,
logit lead and snap ratio under every window; each near tie's other next codes; the start and its
own near ties; $\varphi$; and the minimised machine's blocks. It is stored compressed, with the claim
it supports (the target, the margins, the network file and its hash, and how the network was read as
a machine).

The checker shares no code with the certifier. It reads the target only through its transition and
output tables, and checks two things. The \emph{structure}, without the network: the table is closed
(every next code and branch is a row), $\varphi$ maps the start to the target's start, every entry's
labels and next codes agree with the target under $\varphi$, every lead meets the margin, $\varphi$ is
onto the target's states reachable by whole windows, and the recorded minimised machine is $\varphi$'s kernel (which, for a
homomorphism onto a minimal target, is the coarsest congruence). And \emph{fidelity}: every entry is
run again through the network's forward pass, in shuffled batches that mix codes, so in an order and
batching the certifier never used; the snap is redone from the cells' outputs and the codebooks; the
labels must agree, the snapped code must be one of the table's next codes for the entry, and every
code the forward pass finds within $\rho$ must be among the table's branches.

The trusted base is then the checker, the network's loader, the forward pass and its numerical backend,
the query reading where one is used, the target's tables and the target's minimality (which the kernel
claim needs, and which we checked by Moore's refinement). The certifier's search, its bookkeeping of
$\varphi$, its branching and its minimisation are all outside it. ``Independent'' means sharing no code,
not failing independently: both programs read the same forward pass. And the checker tests the step from
supplied states. It does not run full-sequence inference, so the link from step to sequence rests on the
C2 test above, and, for the earlier study's networks, on their long tests.

\section{Without the target}
\label{sec:target-free}

The certificate needs the target. A discoverer has only the target's answers, and three things can
be checked without it: that the machine read off fits the observations; that the factors' readouts
are the congruences the training said they would be; and that the machine is minimal. None of them
proves equivalence. With a bound $B$ on the target's states, Chow's test set
\citep{chow-1978-testing-software-design-finite-state-machines,vasilevskii-1973-failure-diagnosis-automata}
does: a minimal machine of $n$ states that agrees with a deterministic, resettable black-box target of at
most $B$ states on every word of the set is equivalent to it. The set grows as $|\Sigma|^{B - n + 1}$ in
the slack $B - n$ (Table~\ref{tab:recertify}, last column, at $k$ 1), so the proof is affordable at the
exact bound and goes out of reach as the slack reaches several states. The earlier study measured this: at the exact bound no answers-only stop was
false and every stopped machine was proved; at 1.1 times the bound, of the stops that came early,
Chow's test set refuted every false one it reached and proved none of them
\citep[Appendix F]{barkhuizen-2026-learning-transducers-with-proof}. The target-free proof is sound
where its premise, the bound, holds.

\section{The certificate as an instrument}
\label{sec:instrument}

A decision cheap enough to run at every curriculum step changes how training is run and read. A
run can stop at its first proof, so a certified run is a Las Vegas run: its answer is exact, its time
random, and a rate of success is a cost. The factors' readouts say exactly which congruence each
factor holds. The certificate's failures are informative too: the counts of wrong outputs and wrong
states, the margins and the minimised machine describe a near miss, and in a strict run the first
failing entry says where the network leaves the target. Split ties, counted without the target, measure a network's own
nondeterminism. In \citet{barkhuizen-2026-learning-transducers-with-proof} a run counted as a success
only when certified in both verdicts and passing every long-test case.

It can also act during training, not always for the better. A guard that undid any update making a
right machine wrong kept the proof in 10 of 10 replays of one network, against 1 of 10 without it; in
a training screen it did not help, and it froze two networks' learning by undoing nearly every update
\citep[Appendix B]{barkhuizen-2026-learning-transducers-with-proof}. Stopping at the first proof
avoids the need for it.

\section{Validation}
\label{sec:validation}

Every number measured here is computed from existing networks; no network was trained for this note.

\paragraph{Recertification under the witness.} Of the \pn{inventory-files}{2{,}385} certificate files
on disk, every one recording a certified network was rerun in that verdict, with its table emitted and
checked, along with \pn{linear-uncertified-invalid}{11} uncertified networks as negatives:
\pn{recert-jobs}{1{,}088} certificates of \pn{recert-networks}{685} networks (Table~\ref{tab:recertify}).
Of the \pn{recert-certified-networks}{674} certified networks, \pn{recert-transformer-networks}{404} have
the transformer step and \pn{linear-bound-tables}{270} the linear window of the second family below. The
linear window's networks are successive curriculum steps of the same seeds, not independent draws, and all
the negatives are among its steps. All \pn{recert-certified}{1{,}077} certified tables were valid under the
checker, and the \pn{recert-uncertified}{11} uncertified ones were not, so the checker and the certifier
agreed in all \pn{recert-agree}{1{,}088}. The negatives fail on structure; fidelity's sensitivity to a
corrupted entry is shown by the checker's tests, not by them. All \pn{recert-reproduced}{1{,}088} fresh
certificates reproduced the stored one (its verdict, reachable codes, $\varphi$, wrong counts, near ties
and minimised machine). Re-running every entry of the \pn{recert-entries}{301{,}381{,}245} in shuffled
batches moved \pn{recert-moved}{0} snaps and changed the bits of \pn{recert-not-bitwise}{0} margins. The
largest tables, of \pn{recert-largest-entries}{5{,}314{,}410} entries, took up to
\pn{recert-largest-megabytes}{57}\,MB and \pn{recert-largest-check-minutes}{5} minutes to check.

\begin{table}[t]
\centering
\caption{Recertification under the witness, per family and window size $k$: certificates, reachable
codes, entries ($|Q| \cdot |\Sigma|^k$), the slowest certificate's and check's seconds, the table's
median bytes per entry, and the target-free proof's test set at the exact state bound and at slack 1 and 2
(words; $k$ 1 only; for the first machine, where a family pools several).}
\label{tab:recertify}
\resizebox{\textwidth}{!}{%
\begin{tabular}{lcccrrrrl}
\toprule
family & $k$ & cert. & $|Q|$ & entries & cert.\,s & check\,s & B/entry & Chow's words \\
\midrule
addition & 1 & 357 & 3--14 & 81--378 & 0.0 & 0.0 & 34 & 82 / 2{,}296 / 62{,}074 \\
addition & 2 & 19 & 3--9 & 2{,}187--6{,}561 & 0.3 & 0.3 & 12 & -- \\
addition & 4 & 110 & 3--10 & 1{,}594{,}323--5{,}314{,}410 & 352.4 & 318.9 & 10 & -- \\
division, $p$ 5 & 1 & 198 & 32--137 & 288--1{,}233 & 0.9 & 0.1 & 22 & 904 / 9{,}040 / 82{,}264 \\
division, $p$ 7 & 1 & 46 & 71--242 & 639--2{,}178 & 1.0 & 0.2 & 20 & 1{,}768 / 17{,}680 / 160{,}888 \\
division, $p$ 11 & 1 & 42 & 125--448 & 1{,}125--4{,}032 & 1.8 & 0.5 & 19 & 4{,}360 / 43{,}600 / 396{,}760 \\
division, $p$ 13 & 1 & 48 & 198--456 & 1{,}782--4{,}104 & 2.5 & 0.6 & 18 & 6{,}088 / 60{,}880 / 554{,}008 \\
sum mod 4 & 1 & 28 & 6--9 & 54--81 & 0.0 & 0.0 & 61 & 37 / 370 / 3{,}367 \\
sum mod 5 & 1 & 12 & 7--13 & 63--117 & 0.0 & 0.0 & 48 & 46 / 460 / 4{,}186 \\
move-to-front 3 & 1 & 12 & 12--21 & 48--84 & 0.1 & 0.0 & 61 & 100 / 500 / 2{,}100 \\
move-to-front 4 & 1 & 61 & 37--79 & 185--395 & 0.6 & 0.1 & 28 & 1{,}331 / 7{,}986 / 41{,}261 \\
move-to-front 5 & 1 & 5 & 195--440 & 1{,}170--2{,}640 & 3.1 & 0.2 & 20 & 18{,}746 / 131{,}222 / 806{,}078 \\
planted ($q$ 16) & 1 & 50 & 20--60 & 100--300 & 0.5 & 0.1 & 39 & 324 / 1{,}944 / 10{,}044 \\
random ($q$ 8) & 1 & 6 & 17--19 & 85--95 & 0.1 & 0.0 & 52 & 164 / 984 / 5{,}084 \\
random ($q$ 12) & 1 & 10 & 19--49 & 95--245 & 0.4 & 0.1 & 41 & 305 / 1{,}830 / 9{,}455 \\
random ($q$ 16) & 1 & 8 & 23--33 & 115--165 & 0.4 & 0.1 & 40 & 405 / 2{,}430 / 12{,}555 \\
$S_3$ & 1 & 19 & 10--19 & 40--76 & 0.3 & 0.1 & 72 & 50 / 250 / 1{,}050 \\
$S_4$ & 1 & 51 & 32--85 & 160--425 & 0.6 & 0.1 & 30 & 363 / 2{,}178 / 11{,}253 \\
$S_5$ & 1 & 6 & 172--331 & 1{,}032--1{,}986 & 2.7 & 0.1 & 20 & 2{,}884 / 20{,}188 / 124{,}012 \\
\bottomrule
\end{tabular}}
\end{table}

\paragraph{Faults, and what sampled tests miss.} On one table for each of
\pn{fault-tables}{672} certified networks, single-entry faults were planted: an entry's label moved on,
or its next code moved to one of another target state, at the rarest, median and commonest entries by how
often a uniform word visits them. The certifier caught \pn{faults-caught}{4{,}032} of
\pn{faults}{4{,}032}, with the first failing entry 1 to \pn{fault-longest}{11} windows from the start.
Since every entry is read, a planted fault must be caught; this tests the implementation for bugs, and
the interest is in the comparison that follows.

For each fault we computed exactly, from the Markov chain of the pair (target state, faulty machine's
state) under uniform tokens, the chance that a large sampled test finds it: 100{,}000 uniform words at
128, 256 and 1{,}024 columns, with 1{,}000 more words from each target state's shortest prefix at $k$ 1
(all \pn{faults-wide}{408} faults at $k > 1$ have the uniform words only). This is a model. Addition,
division and the modular sums sample their own strata otherwise, so their rows give this catalogue's
chances, not those of their own long tests.

\pn{faults-likely-missed}{34} faults are more likely missed than found, the worst with chance
\pn{fault-worst-miss}{0.99998} (Table~\ref{tab:faults}): \pn{faults-wide-missed}{28} at addition's
$k$ 4, where $531{,}441$ windows thin out every entry's visits, \pn{faults-wide-median-missed}{2} of them
at median entries; and \pn{faults-division-rarest-missed}{6} of division's
\pn{faults-division-rarest}{330} at its rarest entries, the worst \pn{fault-division-worst}{0.9998}.
Below that line the chances fall quickly: up to \pn{fault-division-below}{0.32} on division's rarest and
\pn{fault-short-worst}{0.11} on addition's rarest at $k$ 2; at most
\pn{fault-elsewhere-worst}{$8 \times 10^{-7}$} at $k$ 1 outside division; and at $k$ 4's commonest entries,
at most \pn{fault-wide-worst-common}{$5 \times 10^{-8}$}. All \pn{faults-rarest-first-window}{1{,}344} of
the \pn{faults-rarest}{1{,}344} rarest faults sit at their row's first window, since a row's windows tie on
visits and ties go to the first. That is the zero window, which the modelled test's final window also
visits, so these counts understate what another window in the row would show. These are planted faults
at entries chosen to be rare: they show where sampled tests are blind, not how often learned networks err
there. That learned errors occur is the earlier study's evidence, the two addition networks of the opening.
The sampled tests cover every target state; they do not cover a network's rarely reached codes and windows.

\begin{table}[t]
\centering
\caption{Planted faults per family: faults, caught by the certificate, the first failing entry's
distance from the start in windows, and the sampled test's worst chance of missing a fault at the
rarest, median and commonest entries, rounded up to one figure below 0.01 ($0$: rounds to zero in
double precision, or found with certainty).}
\label{tab:faults}
\begin{tabular}{lrrcccc}
\toprule
family & faults & caught & windows & rarest & median & commonest \\
\midrule
addition & 2136 & 2136 & 1--4 & 1.00 & 0.51 & $5 \times 10^{-8}$ \\
division, $p$ 5 & 582 & 582 & 1--8 & $6 \times 10^{-5}$ & $0$ & $0$ \\
division, $p$ 7 & 138 & 138 & 1--8 & 0.98 & $0$ & $0$ \\
division, $p$ 11 & 126 & 126 & 1--10 & 1.00 & $0$ & $0$ \\
division, $p$ 13 & 144 & 144 & 1--8 & 0.88 & $0$ & $0$ \\
sum mod 4 & 144 & 144 & 2--4 & $0$ & $0$ & $0$ \\
sum mod 5 & 36 & 36 & 1--4 & $0$ & $0$ & $0$ \\
move-to-front 3 & 36 & 36 & 1--7 & $0$ & $0$ & $0$ \\
move-to-front 4 & 192 & 192 & 1--11 & $2 \times 10^{-74}$ & $0$ & $0$ \\
move-to-front 5 & 24 & 24 & 1--10 & $3 \times 10^{-20}$ & $0$ & $0$ \\
planted ($q$ 16) & 150 & 150 & 1--8 & $3 \times 10^{-99}$ & $0$ & $0$ \\
random ($q$ 8) & 18 & 18 & 2--5 & $0$ & $0$ & $0$ \\
random ($q$ 12) & 30 & 30 & 2--10 & $0$ & $0$ & $0$ \\
random ($q$ 16) & 24 & 24 & 3--4 & $0$ & $0$ & $0$ \\
$S_3$ & 66 & 66 & 1--6 & $0$ & $0$ & $0$ \\
$S_4$ & 162 & 162 & 2--10 & $8 \times 10^{-7}$ & $0$ & $0$ \\
$S_5$ & 24 & 24 & 1--9 & $8 \times 10^{-34}$ & $0$ & $0$ \\
\bottomrule
\end{tabular}
\end{table}

\paragraph{The bounds.} The interval bound ran on one table for each of \pn{bound-networks}{620} of the
\pn{recert-certified-networks}{674} certified networks: the strict table where the network has one, else
its tie-branched table. It left out \pn{bound-left}{54}. Of those, \pn{bound-left-long}{52} are at $k$ 4, of
which only \pn{bound-long}{5} ran, the smallest for each code count, so $k$ 4 is sampled: an interval pass
took at most \pn{bound-long-minutes}{10} minutes, but a table it leaves undecided can take days under the
affine bound. The other two failed their first recertification, and were rerun too late for these passes. Of the
\pn{bound-networks}{620}, \pn{bound-unsupported}{7} use window branches that the implementation does not
cover, leaving \pn{bound-tables}{613} tables (\pn{bound-strict}{369} strict) and
\pn{bound-entries}{8{,}814{,}470} entries. No checked floating-point value (the cells' outputs and, for
intervals, the probabilities) lay outside its enclosure (\pn{bound-outside}{0}) and no recorded label was
contradicted (\pn{bound-contradicted}{0}).

The intervals alone decide \pn{bound-labels-percent}{99.1}\% of labels and \pn{bound-snaps-percent}{98.1}\%
of snaps by entries, but the \pn{bound-long}{5} tables at $k$ 4 dominate that count: on the other
\pn{bound-rest-tables}{608} tables the shares are \pn{bound-rest-labels-percent}{74.1}\% and
\pn{bound-rest-snaps-percent}{45.5}\%. They decide \pn{bound-full}{384} tables wholly, of which
\pn{linear-bound-full}{270} are the linear window's, and \pn{bound-transformer-full}{114} of the
\pn{bound-transformer-tables}{343} with the transformer step (Table~\ref{tab:bound}). On division's deepest
they decide almost no snaps. No snap was ever merely covered, but only \pn{bound-branch-tables}{13} of the
tables have branches at all.

The affine bound then ran on exactly the \pn{affine-tables}{229} tables the intervals left undecided
(\pn{affine-mismatch}{0} differences between the two sets), all \pn{affine-strict}{229} strict,
\pn{affine-entries}{1{,}819{,}953} entries. It decided them all: \pn{affine-labels-undecided}{0} labels
undecided, \pn{affine-labels-contradicted}{0} contradicted, \pn{affine-snaps-uncovered}{0} snaps undecided,
\pn{affine-outside}{0} values outside their forms. The symbol cap was reached in
\pn{affine-capped-tables}{188} tables without costing a decision. The price is time: the longest table took
\pn{affine-longest-hours}{45} hours on one core.

So in all \pn{bounds-decided}{613} tables, the exact real-valued function makes every entry's recorded
decisions, including \pn{bound-near-snaps}{7{,}268} snaps within the near-tie ratio, in
\pn{bound-near-tables}{210} tables, the tightest at ratio \pn{bound-smallest-ratio}{1.0000067}. For these
tables the margins are adequate in fact, under the assumptions of Section~\ref{sec:numerics}, not by
assumption alone. This does not reach the tie-branched tables of networks also certified strict, whose
branches were not bounded, or the start's snap, whose margin is wide: the smallest start ratio is
\pn{bound-start-ratio}{2.53}.

\begin{table}[t]
\centering
\caption{The bounds per family. Intervals: tables, tables wholly decided, the shares (\%) of labels
sound, of snaps sound and of snaps covered. Affine, on the tables the intervals left undecided: tables
and tables wholly decided. Last, checked floating-point values outside their enclosures under either bound.}
\label{tab:bound}
\begin{tabular}{lrrrrrrrr}
\toprule
 & \multicolumn{5}{c}{intervals} & \multicolumn{2}{c}{affine} & \\
family & tables & whole & labels & snaps & covered & tables & whole & outside \\
\midrule
addition & 304 & 296 & 100.00 & 99.97 & 0.00 & 8 & 8 & 0 \\
division, $p$ 5 & 97 & 2 & 99.89 & 37.37 & 0.00 & 95 & 95 & 0 \\
division, $p$ 7 & 23 & 0 & 82.36 & 3.76 & 0.00 & 23 & 23 & 0 \\
division, $p$ 11 & 20 & 0 & 41.60 & 0.28 & 0.00 & 20 & 20 & 0 \\
division, $p$ 13 & 23 & 0 & 19.08 & 0.02 & 0.00 & 23 & 23 & 0 \\
sum mod 4 & 24 & 24 & 100.00 & 100.00 & 0.00 & 0 & 0 & 0 \\
sum mod 5 & 6 & 6 & 100.00 & 100.00 & 0.00 & 0 & 0 & 0 \\
move-to-front 3 & 6 & 6 & 100.00 & 100.00 & 0.00 & 0 & 0 & 0 \\
move-to-front 4 & 29 & 12 & 98.20 & 82.22 & 0.00 & 17 & 17 & 0 \\
move-to-front 5 & 3 & 3 & 100.00 & 100.00 & 0.00 & 0 & 0 & 0 \\
planted ($q$ 16) & 25 & 3 & 94.01 & 60.23 & 0.00 & 22 & 22 & 0 \\
random ($q$ 8) & 3 & 3 & 100.00 & 100.00 & 0.00 & 0 & 0 & 0 \\
random ($q$ 12) & 5 & 5 & 100.00 & 100.00 & 0.00 & 0 & 0 & 0 \\
random ($q$ 16) & 4 & 4 & 100.00 & 100.00 & 0.00 & 0 & 0 & 0 \\
$S_3$ & 11 & 6 & 100.00 & 75.67 & 0.00 & 5 & 5 & 0 \\
$S_4$ & 27 & 11 & 96.33 & 52.86 & 0.00 & 16 & 16 & 0 \\
$S_5$ & 3 & 3 & 100.00 & 100.00 & 0.00 & 0 & 0 & 0 \\
\bottomrule
\end{tabular}
\end{table}

\paragraph{A second step family.} The linear-window networks (the window without its transformer blocks)
are counted above; separately, they number \pn{linear-jobs}{321} certificates of
\pn{linear-networks}{281} networks, \pn{linear-certified}{310} certified and all \pn{linear-valid}{310} of
those valid, and the \pn{linear-uncertified-invalid}{11} others invalid. The interval bound alone decides
all \pn{linear-bound-full}{270} of the \pn{linear-bound-tables}{270} tables it ran on. The method is not tied
to the transformer step, only to the conditions.

\section{Related work}

\paragraph{Extraction.} Quantising a recurrent network's state into a machine is old
\citep{giles-1992-learning-extracting-fsa-second-order-rnn,omlin-1996-extraction-rules-discrete-time-rnn,zeng-1993-self-clustering-recurrent-networks};
active learning with the network as an oracle
\citep{weiss-2018-extracting-automata-queries-counterexamples,mayr-2018-regular-inference-artificial-neural-networks},
register automata \citep{hong-2025-register-automata-data-sequences} and hyperplane partitions
\citep{wickramasinghe-2026-extracting-dfa-hyperplane-partitioning} extend it. The extracted machine is a
surrogate; \citet{omlin-1996-extraction-rules-discrete-time-rnn} found the smallest extracted machine
consistent with the training set generalising better than its network.

\paragraph{Networks that are machines.} Zeng et al.'s self-clustering network is discretised as it runs
and behaves exactly like a state machine; they run each state found under each input, but state no closure
of the search, soundness or decision. Quantised bottlenecks \citep{koul-2019-finite-state-recurrent-policies},
state-regularised networks \citep{wang-2019-state-regularized-rnns} and
\citet{aichernig-2024-learning-minimal-automata-rnn}, who check the machine up to isomorphism, read the
machine off the transitions seen in data, so states no trace visits go unchecked.

The nearest precedent is \citet{brandoit-2026-length-independent-state-tracking}. They read a trained
network's tables, explore every reachable joint index under every symbol, fix a map to the task's states on
first reach, check it, and conclude correctness at every length by induction: the same core as
Section~\ref{sec:procedure}, and earlier. The differences are these. Their step's tables read the symbol
alone, so one pass per symbol gives them; here the step reads the state, so the enumeration takes a forward
pass per code and window. They leave the enumeration over windows in principle for their stacked task; here
it is carried out, to $531{,}441$ windows. Their state is supervised in training, though their check does
not use the labels; here training labels no state. We found in their text no near-tie branching, emitted
witness, independent checker or completeness result, and correctness conditional on a margin they state
without, as far as we read, checking it; those are what this note adds.

\paragraph{Verification.} Neural-network verification checks properties: invariants of recurrent
networks \citep{jacoby-2020-verifying-rnn-invariant-inference}, properties through extracted automata
\citep{khmelnitsky-2021-property-directed-verification-rnn}, and, for quantised networks, bit-exact
properties by SMT, a problem they show PSPACE-hard \citep{henzinger-2021-scalable-verification-quantized-nn}. Our
question, equivalence with a machine at every length, is a property the architecture makes cheap:
the finite interface turns a sequence-level property into a finite enumeration. Floating-point
soundness \citep{jia-2021-exploiting-verified-nn-floating-point} is what Section~\ref{sec:numerics}
addresses. Its tools are not ours: floating-point-sound affine forms for networks are DeepZ's
\citep{singh-2018-fast-effective-robustness-certification}, and affine bounds through attention come from
\citet{shi-2020-robustness-verification-transformers} and DeepT
\citep{bonaert-2021-fast-precise-certification-transformers}. These certify local properties of one forward
pass; here each bound is run per table entry, to check a recorded decision.

\paragraph{Testing and certification.} Conformance testing gives complete test suites for machines of
bounded size \citep{chow-1978-testing-software-design-finite-state-machines,vasilevskii-1973-failure-diagnosis-automata,kocsis-2024-complete-test-suites-monoidal-closed};
the certificate is the white-box counterpart, complete with no bound because it reads the machine's
states directly. The witness follows certifying algorithms \citep{mcconnell-2011-certifying-algorithms}
and translation validation \citep{pnueli-1998-translation-validation}, as do proof-producing network
verifiers \citep{isac-2022-neural-network-verification-proof-production,desmartin-2024-certified-proof-checker-imandra}.

\paragraph{Discrete interfaces.} Snapped codes are vector quantisation
\citep{van-den-oord-2017-vq-vae-neural-discrete-representation}; discrete latent states are common in
world models \citep{hafner-2020-dreamerv2-discrete-world-models}. Where such a latent is the only
channel between steps, is snapped exactly and is evaluated deterministically, the conditions here hold
and the procedure applies. As published, DreamerV2 does not meet them: a continuous recurrent state also
crosses steps, and the latents are sampled. We have not tried the procedure on either.

No part of the method is new on its own; each has the precedent above, and the core decision is
Brandoit et al.'s. What we did not find in our review is the combination for steps that read the state:
stated conditions, the decision with branched ties and a completeness result, a witness with an independent
checker, numerics bounded per decision, and its use as an instrument.

\section{Limits}

\begin{itemize}
\item It applies only to networks built with a finite, exactly snapped interface.
\item The window alphabet grows as $|\Sigma|^k$, so the enumeration does too; a window of four
columns of 27 tokens is \pn{recert-largest-entries}{5{,}314{,}410} entries for
\pn{recert-largest-codes}{10} codes.
\item Without the target, a bound on its states is needed, and the target-free proof grows
exponentially in the bound's slack.
\item The trusted base is the checker, the forward pass and the target's tables.
\item The bounds add an assumption on the maths library and their own unverified programs. They cover
\pn{bound-tables}{613} of \pn{recert-certified-networks}{674} networks, few of them at $k$ 4, and
none with window branches.
\item One architecture family, with a second step; every target is arithmetic, algebraic or a random
machine.
\item A certificate says whether a network is right, never when to stop trying to train one.
\end{itemize}

\section{Reproducibility}

The tables and every number marked in the source are regenerated from the result files by
indrajala-studies' \texttt{scripts/paper/certificate\_note\_tables.py}, whose \texttt{check} compares
them with this text and fails on any difference; the result files are listed with their SHA-256 hashes
in a manifest. Figures quoted from the earlier study are cited to it. The certifier, the checker
(\texttt{scripts/certnote/certificate\_check.py}), the fault injection, the two bounds and the
recertification driver are in the same repository. This note's result files, its tables and the networks
not in the earlier study's datapack \citep{barkhuizen-2026-datapack} are in its addendum,
\ifx\addendumdoi\empty\draftnote{the addendum's DOI}\else\href{https://doi.org/\addendumdoi}{doi:\addendumdoi}\fi.

\section*{Authorship and the use of AI}

David Barkhuizen posed the question this note serves, learning finite-state machines from their
behaviour and known structure, with proof; asked for this note and its framing (methodology, with
results only as validation); and decided what it claims and how it is published. Claude (Opus 5.5,
a model made by Anthropic), working through Claude Code, did the rest: the literature review, the
witness, the checker, the fault injection and the bounds, their runs and analysis, and the drafting.
Before release the note is reviewed against the data and the literature by fresh instances of the
same model, working without the drafting context. Barkhuizen takes responsibility for publishing it,
without having checked every line of its code, which is why the note rests its findings on checks
anyone can rerun. Claude is not listed as an author; Anthropic has not been involved in this work as
an organisation and has not reviewed or endorsed it.

\bibliographystyle{plainnat}
\bibliography{library,extra}

\end{document}
